\documentclass[a4paper,11pt]{article}
\usepackage[a4paper,margin=2cm]{geometry}
\usepackage[T1]{fontenc}
\usepackage[utf8]{inputenc}
\usepackage[ngerman,english]{babel}
\usepackage{lmodern,microtype,parskip,enumitem}
\usepackage[table]{xcolor}
\usepackage[most]{tcolorbox}
\usepackage{titlesec,array,longtable}
\usepackage[hidelinks]{hyperref}
\definecolor{HeaderBlueDark}{RGB}{70,110,170}
\definecolor{RowBlueLight}{RGB}{235,243,255}
\definecolor{RowBlueDark}{RGB}{220,233,250}
\definecolor{SoftGray}{RGB}{90,90,90}
\setlength{\parindent}{0pt}
\setlist[itemize]{leftmargin=1.4em,itemsep=0.35em,topsep=0.35em}
\linespread{1.06}
\renewcommand{\arraystretch}{1.35}
\setlength{\tabcolsep}{5pt}
\titleformat{\section}[block]{\Large\bfseries\color{HeaderBlueDark}}{}{0pt}{}
\titlespacing{\section}{0pt}{1.3em}{0.8em}
\titleformat{\subsection}[block]{\large\bfseries\color{HeaderBlueDark}}{}{0pt}{}
\titlespacing{\subsection}{0pt}{1.0em}{0.6em}
\newtcolorbox{exbox}{enhanced,breakable,colback=RowBlueLight,colframe=RowBlueDark,boxrule=0.4pt,arc=1.5mm,left=1.2mm,right=1.2mm,top=1mm,bottom=1mm,title={Beispiele \textnormal{(Examples)}},fonttitle=\bfseries,colbacktitle=RowBlueDark,coltitle=black}
\NewDocumentEnvironment{grammarentry}{mm+b}{%
\begin{tcolorbox}[enhanced,breakable,colback=white,colframe=HeaderBlueDark,boxrule=0.6pt,arc=1.5mm,left=1.5mm,right=1.5mm,top=1mm,bottom=1mm,colbacktitle=HeaderBlueDark,coltitle=white,fonttitle=\bfseries,title={#1\hfill\normalfont\footnotesize #2}]
#3
\end{tcolorbox}
}{}
\newcommand{\GermanText}[1]{\textbf{Deutsch: }#1\par}
\newcommand{\EnglishText}[1]{{\small\color{SoftGray}\textbf{English: }#1\par}}
\newcommand{\ExampleItem}[2]{\item \textbf{#1}\\{\small\color{SoftGray}#2}}
\newcommand{\EN}[1]{\EnglishText{#1}}
\title{Die vier Kasus\\\large\textnormal{The Four Grammatical Cases}}
\date{}
\begin{document}
    \maketitle
    
    
    \tableofcontents
    \newpage

    \section{Grundlagen: Was ist Kasus? \textnormal{(Basics: what is grammatical case?)}}
        \begin{grammarentry}{Vier grammatische Fälle}{Four cases}
            \GermanText{Der \textbf{Kasus} kennzeichnet die grammatische Rolle und die syntaktischen Beziehungen einer Nominalgruppe. Im Deutschen gibt es \textbf{Nominativ}, \textbf{Akkusativ}, \textbf{Dativ} und \textbf{Genitiv}.}
            \EnglishText{\textbf{Case} marks grammatical roles and syntactic relationships of noun phrases. German has \textbf{nominative}, \textbf{accusative}, \textbf{dative}, and \textbf{genitive}.}
            \GermanText{Kasusformen zeigen sich vor allem an \textbf{Artikeln, Pronomen und Adjektivendungen}; manchmal auch am Nomen. Ein freies Adverb wie \textbf{heute} wird nicht dekliniert.}
            \EnglishText{Case is primarily visible in \textbf{articles, pronouns, and adjective endings}, and sometimes in nouns. An ordinary adverb such as \textbf{heute} is not declined.}
        \end{grammarentry}

        \rowcolors{2}{RowBlueLight}{RowBlueDark}
        \begin{longtable}{>{\bfseries}p{2.7cm}|p{2.5cm}|p{5.1cm}|p{4cm}}
            \rowcolor{HeaderBlueDark}
            \color{white}\textbf{Kasus / Case} &
            \color{white}\textbf{Frage / Question} &
            \color{white}\textbf{Beispiel / Example} &
            \color{white}\textbf{Rolle / Role}\\
            Nominativ & Wer? Was? & \textbf{Der Mann} kommt. & Subjekt / subject\\
            Akkusativ & Wen? Was? & Ich sehe \textbf{den Mann}. & Akkusativobjekt / direct object\\
            Dativ & Wem? & Ich helfe \textbf{dem Mann}. & Dativobjekt / dative object\\
            Genitiv & Wessen? & Das ist das Auto \textbf{des Mannes}. & Attribut / possessive modifier\\
        \end{longtable}

    \section{Subjekt, Objekt und Prädikativ \textnormal{(Subject, object and predicative complement)}}
        \begin{grammarentry}{Drei unterschiedliche Funktionen}{Syntactic functions}
            \GermanText{\textbf{Subjekt}: Das Satzglied, mit dem das finite Verb hinsichtlich Person und Numerus normalerweise übereinstimmt; es steht im \textbf{Nominativ}. \textbf{Objekt}: Eine vom Verb verlangte oder zugelassene Ergänzung, typischerweise im \textbf{Akkusativ}, \textbf{Dativ} oder \textbf{Genitiv} oder als \textbf{Präpositionalobjekt}. \textbf{Prädikativ}: Eine Ergänzung, die Subjekt oder Objekt eine Identität, Rolle oder Eigenschaft zuschreibt.}
            \EnglishText{\textbf{Subject}: the clause element that normally controls finite-verb agreement and takes \textbf{nominative} case. \textbf{Object}: a complement selected or allowed by the verb, typically \textbf{accusative}, \textbf{dative}, \textbf{genitive}, or \textbf{prepositional}. \textbf{Predicative complement}: a phrase attributing identity, role, or property to the subject or an object.}
            \GermanText{Der Ausdruck \textbf{Nominativobjekt} ist in der üblichen deutschen Schulgrammatik \textbf{keine reguläre Objektkategorie}. Nach kopulativem \textbf{sein}, \textbf{werden} und \textbf{bleiben} steht ein nominales Subjektsprädikativ im \textbf{Nominativ}.}
            \EnglishText{\textbf{Nominative object} is \textbf{not a regular object category} in standard German school grammar. After copular \textbf{sein}, \textbf{werden}, and \textbf{bleiben}, a nominal subject complement takes \textbf{nominative} case.}
        \end{grammarentry}

        \begin{exbox}\begin{itemize}
            \ExampleItem{Er ist ein guter Lehrer.}{He is a good teacher. (Er = subject; ein guter Lehrer = nominative subject complement)}
            \ExampleItem{Ich kenne einen guten Lehrer.}{I know a good teacher. (einen guten Lehrer = accusative object)}
            \ExampleItem{Sie nennt ihn einen guten Lehrer.}{She calls him a good teacher. (ihn = accusative object; einen guten Lehrer = accusative object complement)}
            \ExampleItem{Zu spät zu kommen ist ein anderes Problem.}{Being late is another problem. (infinitive clause = subject; ein anderes Problem = nominative complement)}
            \ExampleItem{Die Kinder sind glücklich.}{The children are happy. (glücklich = uninflected adjective, not a nominative noun phrase)}
        \end{itemize}\end{exbox}

    \section{Alle vier Kasus im Detail \textnormal{(All four cases in detail)}}
        \subsection{Nominativ \textnormal{(Nominative)}}
            \begin{grammarentry}{Subjekt und nominales Subjektsprädikativ}{Main roles}
                \GermanText{Der \textbf{Nominativ} bezeichnet das Subjekt und steht auch bei nominalen Subjektsprädikativen nach \textbf{sein}, \textbf{werden} und \textbf{bleiben}. Auch Fragen nach dem Subjekt und dessen Antworten stehen entsprechend im Nominativ.}
                \EnglishText{\textbf{Nominative} marks the subject and nominal subject complements after \textbf{sein}, \textbf{werden}, and \textbf{bleiben}. Questions about the subject and their corresponding answers use the nominative.}
            \end{grammarentry}
            \begin{exbox}\begin{itemize}
                \ExampleItem{Ein neuer Student kommt.}{A new student arrives.}
                \ExampleItem{Der Student wird ein guter Forscher.}{The student becomes a good researcher.}
                \ExampleItem{Wer ist das? Das ist mein Bruder.}{Who is that? That is my brother.}
            \end{itemize}\end{exbox}

        \subsection{Akkusativ \textnormal{(Accusative)}}
            \begin{grammarentry}{Objekt, Präposition und Zeitmaß}{Common roles}
                \GermanText{Der \textbf{Akkusativ} erscheint bei vielen Verbobjekten (\textbf{sehen, kaufen, fragen}), nach bestimmten Präpositionen (\textbf{für, durch, ohne, gegen, um}), bei Richtungsangaben nach Wechselpräpositionen und bei manchen Zeit- und Maßangaben (\textbf{den ganzen Tag}).}
                \EnglishText{The \textbf{accusative} appears with many verb objects (\textbf{see, buy, ask}), after certain prepositions (\textbf{for, through, without, against, around}), with direction after two-way prepositions, and with some expressions of time and measure (\textbf{all day}).}
            \end{grammarentry}
            \begin{exbox}\begin{itemize}
                \ExampleItem{Ich frage den Lehrer.}{I ask the teacher.}
                \ExampleItem{Ich warte auf den Bus.}{I am waiting for the bus. (auf + accusative: governed prepositional object)}
                \ExampleItem{Ich lege das Buch auf den Tisch.}{I put the book onto the table. (direction: accusative)}
                \ExampleItem{Ich arbeite den ganzen Tag.}{I work all day. (adverbial accusative)}
            \end{itemize}\end{exbox}

        \subsection{Dativ \textnormal{(Dative)}}
            \begin{grammarentry}{Dativobjekt und lokale Ergänzung}{Common roles}
                \GermanText{Der \textbf{Dativ} wird von bestimmten Verben (\textbf{helfen, danken, gehören, gefallen}), Präpositionen (\textbf{mit, bei, zu, nach, seit, von, aus}) und Wechselpräpositionen bei einer statischen Ortsangabe gefordert. Er bezeichnet häufig einen Empfänger, ist aber \textbf{nicht immer} ein indirektes Objekt.}
                \EnglishText{The \textbf{dative} is required by certain verbs (\textbf{help, thank, belong to, please}), prepositions (\textbf{with, at, to, after, since, from, out of}), and two-way prepositions for static location. It often marks a recipient but is \textbf{not always} an indirect object.}
            \end{grammarentry}
            \begin{exbox}\begin{itemize}
                \ExampleItem{Ich helfe dem Kind.}{I help the child.}
                \ExampleItem{Sie gibt dem Kind ein Buch.}{She gives the child a book. (dem Kind = dative; ein Buch = accusative)}
                \ExampleItem{Das Buch liegt auf dem Tisch.}{The book is lying on the table. (location: dative)}
            \end{itemize}\end{exbox}

        \subsection{Genitiv \textnormal{(Genitive)}}
            \begin{grammarentry}{Zugehörigkeit und seltene Verbobjekte}{Common roles}
                \GermanText{Der \textbf{Genitiv} steht vor allem bei Nomen als Attribut (\textbf{das Dach des Hauses}), nach formellen Präpositionen wie \textbf{während, trotz, wegen} und bei wenigen Verben wie \textbf{bedürfen}, \textbf{gedenken} oder \textbf{sich einer Sache annehmen}.}
                \EnglishText{The \textbf{genitive} mainly occurs as a noun modifier (\textbf{the roof of the house}), after formal prepositions such as \textbf{during, despite, because of}, and with a few verbs such as \textbf{bedürfen}, \textbf{gedenken}, or \textbf{sich einer Sache annehmen}.}
            \end{grammarentry}
            \begin{exbox}\begin{itemize}
                \ExampleItem{Das ist das Auto meines Bruders.}{That is my brother's car.}
                \ExampleItem{Wegen des Wetters bleiben wir zu Hause.}{Because of the weather, we are staying at home.}
                \ExampleItem{Diese Frage bedarf einer Erklärung.}{This question requires an explanation. (genitive object)}
            \end{itemize}\end{exbox}

    \section{Deklination im Überblick \textnormal{(Overview of declension)}}
        \begin{grammarentry}{Bestimmte Artikel}{Definite articles}
            \GermanText{Die Artikel zeigen besonders deutlich den Kasus. Im Plural wird im Dativ am Nomen häufig ein zusätzliches \textbf{-n} verwendet, sofern die Pluralform nicht bereits auf \textbf{-n} oder \textbf{-s} endet: \textbf{mit den Kindern}.}
            \EnglishText{Articles show case particularly clearly. In the dative plural, nouns often receive an extra \textbf{-n} unless the plural already ends in \textbf{-n} or \textbf{-s}: \textbf{mit den Kindern} (with the children).}
        \end{grammarentry}
        \rowcolors{2}{RowBlueLight}{RowBlueDark}
        \begin{longtable}{>{\bfseries}p{3.2cm}|p{2.8cm}|p{2.8cm}|p{2.8cm}|p{2.6cm}}
            \rowcolor{HeaderBlueDark}
            \color{white}\textbf{Kasus / Case} &
            \color{white}\textbf{Maskulin} &
            \color{white}\textbf{Feminin} &
            \color{white}\textbf{Neutrum} &
            \color{white}\textbf{Plural}\\
            Nominativ & der & die & das & die\\
            Akkusativ & den & die & das & die\\
            Dativ & dem & der & dem & den\\
            Genitiv & des & der & des & der\\
        \end{longtable}
        \begin{grammarentry}{Adjektive und Pronomen}{Agreement and case}
            \GermanText{Innerhalb einer \textbf{Nominalgruppe} richten sich Artikel und Adjektivendung nach Kasus, Genus und Numerus: \textbf{ein guter Lehrer} (Nom.), \textbf{einen guten Lehrer} (Akk.), \textbf{einem guten Lehrer} (Dat.), \textbf{eines guten Lehrers} (Gen.).}
            \EnglishText{Inside a \textbf{noun phrase}, articles and adjective endings depend on case, gender, and number: \textbf{a good teacher} in nominative, accusative, dative, and genitive forms.}
            \GermanText{Personalpronomen verändern sich ebenfalls: \textbf{ich -- mich -- mir -- meiner}; \textbf{er -- ihn -- ihm -- seiner}. Die prädikativen Adjektive dagegen bleiben normalerweise unverändert: \textbf{Der Lehrer ist gut}.}
            \EnglishText{Personal pronouns also change form: \textbf{ich -- mich -- mir -- meiner}; \textbf{er -- ihn -- ihm -- seiner}. Predicative adjectives, in contrast, normally remain unchanged: \textbf{The teacher is good}.}
        \end{grammarentry}

    \section{Kasus wird durch Konstruktionen bestimmt \textnormal{(Case depends on constructions)}}
        \begin{grammarentry}{Verb, Präposition, Adjektiv und Nomen}{Case assignment}
            \GermanText{\textbf{1. Verb}: \textbf{jemandem helfen} (Dat.), \textbf{jemanden sehen} (Akk.). \textbf{2. Präposition}: \textbf{mit jemandem} (Dat.), \textbf{für jemanden} (Akk.). \textbf{3. Adjektiv}: \textbf{jemandem dankbar sein} (Dat.). \textbf{4. Nomen}: \textbf{die Erinnerung an jemanden} (\textbf{an} + Akk.). \textbf{5. Satzfunktion}: Subjekt und nominales Subjektsprädikativ (Nom.).}
            \EnglishText{\textbf{1. Verb}: help someone (dative), see someone (accusative). \textbf{2. Preposition}: with someone (dative), for someone (accusative). \textbf{3. Adjective}: grateful to someone (dative). \textbf{4. Noun}: memory of someone (\textbf{an} + accusative). \textbf{5. Clause function}: subject and nominal subject complement (nominative).}
            \GermanText{Die \textbf{Wortstellung} allein entscheidet nicht über den Kasus: \textbf{Den Mann sieht die Frau} hat dieselben Kasusrollen wie \textbf{Die Frau sieht den Mann}.}
            \EnglishText{\textbf{Word order} alone does not determine case: \textbf{Den Mann sieht die Frau} and \textbf{Die Frau sieht den Mann} contain the same accusative object and nominative subject.}
        \end{grammarentry}


    
    
    \section{Arten von Objekten \textnormal{(Types of objects)}}
        \begin{grammarentry}{Nicht jedes Objekt ist ein Akkusativobjekt}{Object classification}
            \GermanText{Die gewöhnlichen Objektarten heißen \textbf{Akkusativobjekt}, \textbf{Dativobjekt}, \textbf{Genitivobjekt} und \textbf{Präpositionalobjekt}. Die Fragewörter \textbf{wen/was}, \textbf{wem}, \textbf{wessen} helfen, aber die \textbf{Rektion des Verbs} ist entscheidend. Ein \textbf{Nominativobjekt} zählt in der traditionellen Schulgrammatik nicht zu diesen Objektarten.}
            \EnglishText{The standard object types are \textbf{accusative}, \textbf{dative}, \textbf{genitive}, and \textbf{prepositional objects}. The question words \textbf{whom/what}, \textbf{to whom}, and \textbf{whose} can help, but the \textbf{verb's government pattern} is decisive. A \textbf{nominative object} is not a standard object category in traditional school grammar.}
        \end{grammarentry}
        \begin{exbox}\begin{itemize}
            \ExampleItem{Ich besuche meinen Freund.}{I visit my friend. (meinen Freund = accusative object)}
            \ExampleItem{Ich vertraue meinem Freund.}{I trust my friend. (meinem Freund = dative object)}
            \ExampleItem{Wir gedenken der Opfer.}{We commemorate the victims. (der Opfer = genitive object)}
            \ExampleItem{Wir warten auf den Zug.}{We are waiting for the train. (auf den Zug = prepositional object; auf + accusative)}
            \ExampleItem{Sie spricht mit ihrem Bruder über das Problem.}{She speaks with her brother about the problem. (mit + dative; über + accusative)}
        \end{itemize}\end{exbox}
        \begin{grammarentry}{Präpositionalobjekt oder freie Angabe?}{Prepositional object or adjunct?}
            \GermanText{Ein \textbf{Präpositionalobjekt} ist eng an die Bedeutung/Rektion eines Verbs gebunden, z.\,B. \textbf{warten auf + Akkusativ}. Eine \textbf{freie adverbiale Angabe} ergänzt dagegen z.\,B. Ort, Zeit oder Art und Weise: \textbf{Ich lese im Garten}. Der Kasus innerhalb einer Präpositionalgruppe wird durch deren Präposition bestimmt.}
            \EnglishText{A \textbf{prepositional object} is governed by the verb, e.g. \textbf{warten auf + accusative}. A \textbf{free adverbial adjunct} provides information such as place, time or manner: \textbf{Ich lese im Garten} (I read in the garden). The case inside a prepositional phrase is determined by its preposition.}
        \end{grammarentry}

    \section{Wechselpräpositionen \textnormal{(Two-way prepositions)}}
        \begin{grammarentry}{Ort und Richtung}{Location and direction}
            \GermanText{Die neun häufigsten Wechselpräpositionen sind \textbf{an, auf, hinter, in, neben, über, unter, vor, zwischen}. Bei einer statischen Lage steht meist \textbf{Dativ} (Wo?), bei einer Zielrichtung/Ortsveränderung \textbf{Akkusativ} (Wohin?). Entscheidend ist \textbf{Ort gegenüber Ziel}, nicht bloß, ob Bewegung stattfindet.}
            \EnglishText{The nine common two-way prepositions are \textbf{an, auf, hinter, in, neben, über, unter, vor, zwischen}. Static location normally takes the \textbf{dative} (where?); destination/change of location takes the \textbf{accusative} (where to?). The key is \textbf{location versus destination}, not merely whether there is movement.}
        \end{grammarentry}
        \begin{exbox}\begin{itemize}
            \ExampleItem{Ich laufe im Park.}{I run in the park. (where? in + dative, despite the movement)}
            \ExampleItem{Ich laufe in den Park.}{I run into the park. (where to? in + accusative)}
            \ExampleItem{Das Bild hängt an der Wand.}{The picture hangs on the wall. (location, dative)}
            \ExampleItem{Ich hänge das Bild an die Wand.}{I hang the picture on the wall. (destination, accusative)}
        \end{itemize}\end{exbox}

    \section{Apposition und Kasus \textnormal{(Apposition and case)}}
        \begin{grammarentry}{Gleichsetzung mit einem Nomen}{Apposition}
            \GermanText{Eine \textbf{Apposition} erläutert eine Nominalgruppe. Sie übernimmt typischerweise deren Kasus: \textbf{Mein Bruder, ein Arzt, wohnt hier} (Nominativ) und \textbf{Ich spreche mit meinem Bruder, einem Arzt} (Dativ). Eine Apposition ist von einem \textbf{Prädikativ nach sein} zu unterscheiden.}
            \EnglishText{An \textbf{apposition} explains a noun phrase and typically shares its case: \textbf{Mein Bruder, ein Arzt, wohnt hier} (nominative) and \textbf{Ich spreche mit meinem Bruder, einem Arzt} (dative). An apposition differs from a \textbf{predicative complement after sein}.}
        \end{grammarentry}

    \section{Kasus bei deutschen Verben \textnormal{(Case with German verbs)}}
        
        \textbf{Im Deutschen reicht es nicht immer, nur direktes und indirektes Objekt zu erkennen.}
        \EN{In German, it is not always enough to recognize only direct and indirect objects.}
        
        \textbf{Oft bestimmt das Verb selbst, welchen Kasus das Objekt bekommt.}
        \EN{Often the verb itself determines which case the object takes.}
        
        \textbf{Deshalb sollte man viele Verben zusammen mit ihrem Kasus lernen.}
        \EN{Therefore, many verbs should be learned together with their case.}
        
        \subsection{Verben mit Dativ \textnormal{(Verbs with dative)}}
            
            \textbf{Einige Verben brauchen immer oder meistens den Dativ.}
            \EN{Some verbs always or usually require the dative.}
            
            \textbf{helfen + Dativ}
            \EN{to help + dative}
            
            \textbf{Ich helfe dir.}
            \EN{I help you.}
            
            \textbf{danken + Dativ}
            \EN{to thank + dative}
            
            \textbf{Ich danke dem Großvater.}
            \EN{I thank the grandfather.}
            
            \textbf{antworten + Dativ}
            \EN{to answer + dative}
            
            \textbf{Ich antworte der Lehrerin.}
            \EN{I answer the teacher.}
            
            \textbf{gehören + Dativ}
            \EN{to belong to + dative}
            
            \textbf{Das Buch gehört mir.}
            \EN{The book belongs to me.}
        
        \subsection{Verben mit Akkusativ \textnormal{(Verbs with accusative)}}
            
            \textbf{Die meisten Verben mit Objekt brauchen den Akkusativ.}
            \EN{Most verbs with an object require the accusative.}
            
            \textbf{sehen + Akkusativ}
            \EN{to see + accusative}
            
            \textbf{Ich sehe dich.}
            \EN{I see you.}
            
            \textbf{haben + Akkusativ}
            \EN{to have + accusative}
            
            \textbf{Ich habe einen Kaffee.}
            \EN{I have a coffee.}
            
            \textbf{kaufen + Akkusativ}
            \EN{to buy + accusative}
            
            \textbf{Ich kaufe das Brot.}
            \EN{I buy the bread.}
        
        \subsection{Verben mit zwei Objekten \textnormal{(Verbs with two objects)}}
            
            \textbf{Wenn ein Satz zwei Objekte hat, ist die Person oft Dativ und die Sache Akkusativ.}
            \EN{When a sentence has two objects, the person is often dative and the thing is accusative.}
            
            \textbf{Ich gebe dir das Buch.}
            \EN{I give you the book.}
            
            \textbf{dir = Dativ}
            \EN{to you = dative}
            
            \textbf{das Buch = Akkusativ}
            \EN{the book = accusative}
            
            \textbf{Ich werfe dir den Ball zu.}
            \EN{I throw the ball to you.}
            
            \textbf{dir = Dativ}
            \EN{to you / receiver = dative}
            
            \textbf{den Ball = Akkusativ}
            \EN{the ball = accusative}
        
        \subsection{Wichtige Regel \textnormal{(Important rule)}}
            
            \textbf{Für Deutsch muss man zuerst lernen, welchen Kasus ein Verb verlangt.}
            \EN{For German, one must first learn which case a verb requires.}
            
            \textbf{Direktes und indirektes Objekt helfen beim Verstehen, aber sie reichen nicht immer aus.}
            \EN{Direct and indirect objects help with understanding, but they are not always enough.}
            
            \textbf{Die beste Methode ist: Verb + Kasus + Beispiel lernen.}
            \EN{The best method is: learn verb + case + example.}
            
            \textbf{helfen + Dativ: Ich helfe dir.}
            \EN{to help + dative: I help you.}
            
            \textbf{sehen + Akkusativ: Ich sehe dich.}
            \EN{to see + accusative: I see you.}
            
            \textbf{geben + Dativ + Akkusativ: Ich gebe dir das Buch.}
            \EN{to give + dative + accusative: I give you the book.}
    
    
    \section{Kasus bei Präpositionen \textnormal{(Case with prepositions)}}
        
        \textbf{Nicht nur Verben bestimmen den Kasus. Auch Präpositionen bestimmen oft den Kasus.}
        \EN{Not only verbs determine the case. Prepositions often determine the case too.}
        
        \textbf{Eine Präposition ist ein kleines Wort wie mit, für, seit, wegen oder in.}
        \EN{A preposition is a small word like with, for, since, because of, or in.}
        
        \textbf{Nach manchen Präpositionen kommt immer Dativ.}
        \EN{After some prepositions, the dative always comes.}
        
        \textbf{Nach manchen Präpositionen kommt immer Akkusativ.}
        \EN{After some prepositions, the accusative always comes.}
        
        \textbf{Manche Präpositionen können Dativ oder Akkusativ haben.}
        \EN{Some prepositions can take dative or accusative.}
        
        \subsection{Präpositionen mit Dativ \textnormal{(Dative prepositions)}}
            
            \textbf{Diese Präpositionen brauchen immer den Dativ: aus, bei, mit, nach, seit, von, zu.}
            \EN{These prepositions always require the dative: from/out of, at/near, with, after/to, since/for, from/of, to.}
            
            \textbf{mit + Dativ}
            \EN{with + dative}
            
            \textbf{Ich spreche mit dem Lehrer.}
            \EN{I speak with the teacher.}
            
            \textbf{bei + Dativ}
            \EN{at / near / with + dative}
            
            \textbf{Ich bin bei meiner Mutter.}
            \EN{I am at my mother's place.}
            
            \textbf{nach + Dativ}
            \EN{after / to + dative}
            
            \textbf{Nach dem Kurs gehe ich nach Hause.}
            \EN{After the course I go home.}
            
            \textbf{von + Dativ}
            \EN{from / of + dative}
            
            \textbf{Das ist ein Brief von meinem Bruder.}
            \EN{That is a letter from my brother.}
            
            \textbf{zu + Dativ}
            \EN{to + dative}
            
            \textbf{Ich gehe zu meinem Professor.}
            \EN{I go to my professor.}
        
        \subsection{Die Präposition seit \textnormal{(The preposition seit)}}
            
            \textbf{Seit ist eine Präposition mit Dativ.}
            \EN{Seit is a preposition with the dative.}
            
            \textbf{Man benutzt seit für eine Zeit, die in der Vergangenheit begonnen hat und bis jetzt dauert.}
            \EN{Seit is used for a time that began in the past and continues until now.}
            
            \textbf{seit + Dativ}
            \EN{since / for + dative}
            
            \textbf{Ich lerne seit einem Jahr Deutsch.}
            \EN{I have been learning German for one year.}
            
            \textbf{Er wohnt seit drei Monaten in Wien.}
            \EN{He has been living in Vienna for three months.}
            
            \textbf{Wir kennen uns seit letzter Woche.}
            \EN{We have known each other since last week.}
            
            \textbf{Ich warte seit zehn Minuten.}
            \EN{I have been waiting for ten minutes.}
            
            \textbf{Ich lebe seit 2024 in Österreich.}
            \EN{I have been living in Austria since 2024.}
        
        \subsection{Wichtige Formen mit seit \textnormal{(Important forms with seit)}}
            
            \textbf{seit einem Tag}
            \EN{for one day}
            
            \textbf{seit einer Woche}
            \EN{for one week}
            
            \textbf{seit einem Monat}
            \EN{for one month}
            
            \textbf{seit einem Jahr}
            \EN{for one year}
            
            \textbf{seit zwei Tagen}
            \EN{for two days}
            
            \textbf{seit drei Wochen}
            \EN{for three weeks}
            
            \textbf{seit vier Monaten}
            \EN{for four months}
            
            \textbf{seit fünf Jahren}
            \EN{for five years}
            
            \textbf{seit dem Morgen}
            \EN{since the morning}
            
            \textbf{seit der Kindheit}
            \EN{since childhood}
        
        \subsection{Seit mit Präsens \textnormal{(Seit with present tense)}}
            
            \textbf{Im Deutschen benutzt man mit seit oft das Präsens, wenn die Situation noch jetzt gilt.}
            \EN{In German, with seit one often uses the present tense when the situation is still true now.}
            
            \textbf{Ich wohne seit einem Jahr in Paris.}
            \EN{I have been living in Paris for one year.}
            
            \textbf{Ich arbeite seit zwei Stunden.}
            \EN{I have been working for two hours.}
            
            \textbf{Sie wartet seit zehn Minuten.}
            \EN{She has been waiting for ten minutes.}
            
            \textbf{Wir lernen seit drei Monaten Deutsch.}
            \EN{We have been learning German for three months.}
        
        \subsection{Häufiger Fehler mit seit \textnormal{(Common mistake with seit)}}
            
            \textbf{Falsch: Ich freue mich, dass du hast seit ein Jahr in Paris leben.}
            \EN{Wrong: I am happy that you have since one year live in Paris.}
            
            \textbf{Richtig: Ich freue mich, dass du seit einem Jahr in Paris lebst.}
            \EN{Correct: I am happy that you have been living in Paris for one year.}
            
            \textbf{Nach dass steht das konjugierte Verb am Ende.}
            \EN{After dass, the conjugated verb goes to the end.}
            
            \textbf{Deshalb sagt man: dass du ... lebst.}
            \EN{Therefore one says: that you ... live.}
            
            \textbf{Nach seit kommt Dativ.}
            \EN{After seit comes dative.}
            
            \textbf{Deshalb sagt man: seit einem Jahr.}
            \EN{Therefore one says: for one year.}
            
            \textbf{Mit seit benutzt man oft Präsens, wenn die Handlung noch dauert.}
            \EN{With seit one often uses the present tense when the action is still continuing.}
        
        \subsection{Seit und seid \textnormal{(Seit and seid)}}
            
            \textbf{Seit und seid klingen gleich, aber sie haben verschiedene Bedeutungen.}
            \EN{Seit and seid sound the same, but they have different meanings.}
            
            \textbf{seit = Präposition für Zeit}
            \EN{seit = preposition for time}
            
            \textbf{Ich lerne seit einem Jahr Deutsch.}
            \EN{I have been learning German for one year.}
            
            \textbf{seid = Verbform von sein}
            \EN{seid = verb form of to be}
            
            \textbf{Ihr seid sehr freundlich.}
            \EN{You are very friendly.}
            
            \textbf{Ihr seid seit einer Stunde hier.}
            \EN{You have been here for one hour.}
    
    
    \section{Zusammenfassung \textnormal{(Summary)}}
        
        \textbf{Bei deutschen Verben muss man oft lernen, welchen Kasus das Verb verlangt.}
        \EN{With German verbs, one often has to learn which case the verb requires.}
        
        \textbf{Bei deutschen Präpositionen muss man ebenfalls lernen, welchen Kasus die Präposition verlangt.}
        \EN{With German prepositions, one also has to learn which case the preposition requires.}
        
        \textbf{helfen + Dativ: Ich helfe dir.}
        \EN{to help + dative: I help you.}
        
        \textbf{sehen + Akkusativ: Ich sehe dich.}
        \EN{to see + accusative: I see you.}
        
        \textbf{geben + Dativ + Akkusativ: Ich gebe dir das Buch.}
        \EN{to give + dative + accusative: I give you the book.}
        
        \textbf{seit + Dativ: Ich lerne seit einem Jahr Deutsch.}
        \EN{since / for + dative: I have been learning German for one year.}
        
        \textbf{Die beste Methode ist: Wort + Kasus + Beispiel lernen.}
        \EN{The best method is: learn word + case + example.}

%        \nocite{*}
%        \bibliographystyle{plainnat}
%        \bibliography{/home/donkarlo/Dropbox/repo/nd_shared_research/bibliography/bibliography.bib}

\end{document}